% ID: SEL-D0029-OB-E04
% ESTADO: cotejada. Segunda pasada visual completa y borrador compilado revisados.
% Fuente: Selectividad/matematicas_aplicadas/suplente septiembre/Examen-Matemáticas Aplicadas a las Ciencias Sociales.pdf
% SHA-256: 4b8884e307df4f5ddaf079e0f10be3d3f74d4793874eb03b355d34f4856579cb
% Página: 2; opción B; ejercicio original 4.

Una cadena de supermercados desea estimar la proporción de clientes que adquiere un determinado producto. Para ello ha tomado una muestra aleatoria simple de 1000 clientes y ha observado que 300 compraban ese producto.
\begin{enumerate}[label=\alph*)]
\item \textbf{(1.5 puntos)} Halle, con un nivel de confianza del 95\%, un intervalo de confianza para estimar la proporción de clientes del supermercado que compra ese producto.
\item \textbf{(1 punto)} Si en otra muestra la proporción de clientes que compra ese producto es de 0.25 y el error cometido en la estimación ha sido inferior a 0.03, con un nivel de confianza del 92.5\%, calcule el tamaño mínimo de la muestra.
\end{enumerate}
